% Class-compatibility preamble retained for a reader excerpt of
%   Sahar Diskin, Anna Geisler, "Perfect Matching in Product Graphs and in their
%   Random Subgraphs", arXiv:2404.14020v4, CC BY 4.0.
%
% The source pins its own publisher class (scrartcl); the delivery wrapper accepts
% only the classes article and amsart, so this file supplies the amsart baseline and
% then copies the author's own declaration block mechanically and unchanged
% (pinned-original-source.tex, source lines 30-57: the author's newtheorem
% declarations and the author's macro declarations).  The author's own package
% list of source lines 4-21 is reused only for the entries the retained text needs
% (enumitem, amsmath, amssymb, amsthm, mathtools, xcolor, geometry) and the
% cross-reference packages already needed by the author's own \Cref/\eqref usage
% (hyperref, cleveref); tikz, caption, subcaption, float, graphicx, wrapfig and
% abstract serve figures and the title block that this excerpt does not carry.
% No mathematical macro, statement, symbol or assumption is introduced, changed
% or removed by this file.
\documentclass[10pt,a4paper]{amsart}
\usepackage[shortlabels]{enumitem}
\usepackage{amsmath}
\usepackage{amssymb}
\usepackage{amsthm}
\usepackage{mathtools}
\usepackage{xcolor}
\usepackage[margin=1in]{geometry}
\usepackage{hyperref}
\usepackage{cleveref}
\hypersetup{
  colorlinks=true,
  linkcolor=blue,
  filecolor=magenta,
  urlcolor=cyan
}

\newtheorem{theorem}{Theorem}
\newtheorem{thm}{Theorem}[section]
\newtheorem{corollary}[thm]{Corollary}
\newtheorem{lemma}[thm]{Lemma}
\newtheorem{proposition}[thm]{Proposition}
\newtheorem{claim}[thm]{Claim}
\newtheorem{remark}[thm]{Remark}
\newtheorem{question}[thm]{Question}

\newcommand{\whp}{whp}
\newcommand{\defi}{\operatorname{def}}
\newcommand{\odd}{\operatorname{odd}}
\newcommand{\iso}{\operatorname{iso}}
\newcommand{\cN}{\mathcal{N}}
\newcommand{\cO}{\mathcal{O}}
\newcommand{\cE}{\mathcal{E}}
\newcommand{\cK}{\mathcal{K}}
\newcommand{\cP}{\mathcal{P}}
\newcommand{\cQ}{\mathcal{Q}}

\newcommand{\footremember}[2]{%
  \footnote{#2}%
  \newcounter{#1}%
  \setcounter{#1}{\value{footnote}}%
}

\newcommand{\ag}[1]{\textcolor{purple}{(Anna: #1)}}
\newcommand{\sd}[1]{\textcolor{red}{(Sahar: #1)}}
\title{An exact edge-counting identity for the Gallai-Edmonds parameters of a graph}
\author{Sahar Diskin, Anna Geisler}
\date{arxiv:2404.14020v4; CC BY 4.0}
\begin{document}
\maketitle
\noindent\textit{Editorial note.} All delivered mathematical text is the author's own, copied from the pinned author source sources/combinatorics-fl04-vspace-1-5cm-perfect-matching-in-p/pinned-original-source.tex; no statement or proof step was rewritten, reordered or removed.\par\medskip
\noindent\textit{Editorial note.} Preamble provenance: the source pins the publisher class scrartcl, which the delivery wrapper does not accept (supported wrapper classes are article and amsart), so this card uses sources/combinatorics-fl04-vspace-1-5cm-perfect-matching-in-p/retained-preamble.tex, which supplies the amsart baseline and then carries the author's own newtheorem and macro declarations (pinned source lines 30-57) verbatim. No macro was invented and no mathematical symbol was changed.\par\medskip
\noindent\textit{Editorial note.} Bibliography provenance: the source itself carries no thebibliography environment and ends with bibliographystyle abbrv plus bibliography perc (author source lines 1691-1692), so the bibliography is resolved with BibTeX; the database used is the identical copy staged with the same arXiv version under sources/product-graph-matchings-2404.14020v4/perc.bib, whose licence evidence is also status supported for source\_id arxiv:2404.14020v4. No entry was edited and the bibliography style is the author's own abbrv.\par\medskip
\noindent\textit{Editorial note.} The counted statement is itself the author's displayed identity labelled eq:ge-exact at source line 1100, delivered verbatim inside the statement span, so its label resolves inside the assembled material and no stub is needed for it.\par\medskip
\noindent\textit{Editorial note.} Two displays inside the counted proof carry the author's labels eq:summation-D and eq:summation-A and are combined by the proof itself; they are inside the delivered span and are never cross-referenced from outside it.\par\medskip
\noindent\textit{Editorial note.} The counted statement and proof cite nothing, so the card needs no bibliography entry beyond the BibTeX database declared for the source.\par\medskip
\noindent\textit{Allowed standard prerequisites (editorial).}\begin{itemize}
\item The Gallai-Edmonds decomposition (D, A, C) of a graph, its parameters r and h, and the components D\_i of G[D]. Set up in the author's subsection on the Edmonds-Gallai decomposition (source lines 1054-1098), where the parameters r and h are defined. Fixes every symbol that appears in the identity, in particular the sets A, C, D, their sizes a, r, and the quantity h.
\end{itemize}\medskip
\begin{lemma}\label{lem:ge-boundary-identity}
\begin{equation}\label{eq:ge-exact}
 2h=dr+2e_G(A)+e_G(A\cup D,C)  +\sum_{i:|D_i|\ge3}\bigl(e_G(D_i,D_i^c)-d\bigr).
\end{equation}
\end{lemma}

\begin{proof}
Degree summation in $D$ gives
\[
    d |D|= 2 \sum_{i} e_G(D_i)+ 2 \sum_{i<j} e_G(D_i, D_j) + e_G(D, A) + e_G(D, C).
\]
Rearranging for $h=\sum_{i<j}e_G(D_i,D_j)+e_G(D,C)$, we obtain
\begin{align} \label{eq:summation-D}
    2h=d|D|-2\sum_ie_G(D_i)-e_G(D,A)+e_G(D,C).
\end{align}
Similarly, degree summation in $A$ gives
\begin{align} \label{eq:summation-A}
    d |A| = 2 e_G(A) + e_G(A, D) + e_G(A, C).
\end{align}
Subtracting the LHS and adding the RHS of \eqref{eq:summation-A} to \eqref{eq:summation-D}, we obtain
\[
2h= d( |D|-|A|) - 2 \sum_{i} e_G(D_i) + 2 e_G(A) + e_G(A, C) + e_G(D, C).
\]
Now
\[
 |D|-|A|=r+\sum_{|D_i|\ge3}(|D_i|-1),  \qquad 2e_G(D_i)=d|D_i|-e_G(D_i,D_i^c),
\]
and since each $|D_i|$ is odd both sums over $i$ may be combined, proving the claim.
\end{proof}

\bibliographystyle{plain}
\bibliography{perc}
\end{document}
